% !TeX program = xelatex
\documentclass[a4paper,10pt, twocolumn]{article}
\usepackage[a4paper,top=2cm,bottom=3cm,left=2cm,right=2cm,footskip=0.8cm]{geometry}
\usepackage{iftex}
\ifPDFTeX
  \usepackage[T1]{fontenc}
  \DeclareUnicodeCharacter{2212}{\ensuremath{-}}
  \usepackage{mathptmx,helvet}
\else
  \usepackage{fontspec}
  \IfFontExistsTF{Times New Roman}{\setmainfont{Times New Roman}}{\setmainfont{TeX Gyre Termes}}
  \IfFontExistsTF{Arial}{\setsansfont{Arial}}{\setsansfont{TeX Gyre Heros}}
\fi
\usepackage[brazilian]{babel}
\usepackage{csquotes,microtype,tikz,booktabs}
\usepackage[hidelinks]{hyperref}
\usepackage{graphicx}
\usepackage{float}
\usepackage{caption}    % Para personalizar legendas, opcional
\usepackage{xcolor}
\usepackage{fancyhdr}
\usepackage{titlesec}     % Para personalizar \section
\usepackage[style=ieee, sorting=none,maxbibnames=3,minbibnames=1, language=portuguese]{biblatex}
\definecolor{grayFEI}{HTML}{7F7F7F}
\usepackage{indentfirst}  % Indenta o primeiro parágrafo
\setlength{\parindent}{0.5cm}  % Aumenta a indentação para 2em (pode ser ajustado
\titlespacing*{\section}{0pt}{*0}{0pt}
\usepackage{ragged2e}


%-------------------------- ESTILOS NÃO ALTERAR ------------------------------------------------%
% Estilo para legendas de figura e tabela
\captionsetup[figure]{
  labelfont={color=azulFEI, bf, small, sf}, % "Figura" azul, negrito, Arial-like
  textfont={small, sf},                    % Texto da legenda em Arial-like
  font=footnotesize,                       % Tamanho aproximado de 9pt
  labelsep=period,                         % "Figura 1." e não "Figura 1:"
  name=Figura,                             % Tradução para "Figura"
  justification=raggedright,              % Alinhado à esquerda
  singlelinecheck=false                   % Garante alinhamento mesmo em uma linha
}

\captionsetup[table]{
  labelfont={color=azulFEI, bf, small, sf},
  textfont={small, sf},
  font=footnotesize,
  labelsep=period,
  name=Tabela,
  justification=raggedright,
  singlelinecheck=false
}

% Formatação section e subsection
\titleformat{\subsection}
  {\color{grayFEI}\normalfont\sffamily\selectfont\itshape\fontsize{10}{12}\selectfont}
  {\thesubsection.}{0.2em}{}
\renewcommand{\thesubsection}{\Alph{subsection}}

\titleformat{\section}
  {\color{azulFEI}\normalfont\sffamily\selectfont\bfseries\fontsize{11}{13}\selectfont} % Fonte Helvetica, tamanho 11pt, negrito
  {\thesection.}{0.2em}{}
\renewcommand{\thesection}{\Roman{section}} 
\titlespacing*{\section}
{0pt}{10pt}{0pt}

%Formatação Título
\newcommand{\inovatitle}[1]{%
  \noindent{\sffamily\bfseries\fontsize{18}{20}\selectfont\textcolor{azulFEI}{#1}\par}\vspace{12pt}
}

%Formatação Autores
\newcommand{\autores}[1]{
  \helvet{\fontsize{11}{13}\selectfont\textbf{\textit{#1}}}
}
\newcommand{\orientador}[1]{
  \helvet{\fontsize{11}{13}\selectfont\textbf{\textit{#1}}}
}
\newcommand{\filiacao}[1]{
  \helvet{\fontsize{10}{12}\selectfont{\textit{#1}}}
}
\newcommand{\email}[1]{
  \helvet{\fontsize{10}{12}\selectfont{\textit{#1}}}
}
\newcommand{\resumo}[1]{%
\vspace{15pt}
  \noindent
  {\helvet{\fontsize{12}{20}\selectfont\textbf{\textcolor{azulFEI}{Resumo:}}}} %
  \justifying #1\par\smallskip %
  \par
}

\newcommand{\keyword}[1]{
  \noindent\helvet{\fontsize{12}{20}\selectfont\textbf{\textcolor{azulFEI}{Palavras-chave:}}} #1  \par\vspace{12pt}
}

\renewcommand{\thetable}{\Roman{table}}


% Define a cor azul específica (exemplo: azul FEI)
\definecolor{azulFEI}{HTML}{0070C0} % substitua pelo código que quiser

%Formatação Cabeçalho e rodapé
\setlength{\headsep}{1cm}
\newcommand{\helvet}[1]{{\sffamily\selectfont #1}}
\pagestyle{fancy}
\fancyhead{} % limpa cabeçalho
\renewcommand{\headrulewidth}{0.3pt} % ← ATIVA a linha (espessura ajustável)
\setlength\headheight{1.2cm}
\lhead{\includegraphics[width=4.97cm, height=1cm]{logoinova.png}}
\newcommand{\semester}{2º semestre/2026}
\rhead{\helvet{\semester}}

\fancyfoot{} % limpa rodapé
\lfoot{
  \sffamily\fontsize{8}{10}\selectfont
  Anais dos Trabalhos de Conclusão de Curso do Centro Universitário FEI, São Bernardo do Campo – SP\\
  \vspace{-4pt}\rule{\textwidth}{1pt} % ← linha abaixo do texto do rodapé
}
\rfoot{\sffamily\fontsize{8}{10}\selectfont \thepage}

%-------------------------------- FIM ESTILOS ---------------------------------------------------%



%Adicionar a bibliografia
\addbibresource{referencias.bib}


\begin{document}
\setlength{\columnsep}{20pt}
\raggedbottom
\twocolumn[
\begin{@twocolumnfalse}
\raggedright
\inovatitle{Um \textit{Harness} de IA para Resolução de Perguntas em Linguagem Natural: Adaptando o \textit{Ralph Wiggum Loop} Além do Desenvolvimento de \textit{Software}}

\autores{Hugo Emílio Nomura, Pedro Henrique Correia de Oliveira,\\
Pedro Henrique Satoru Lima Takahashi, Vitor Monteiro Vianna}\par
\filiacao{Ciência da Computação -- Centro Universitário FEI, \textit{campus} São Paulo}\par
\email{unifhnomura@fei.edu.br, unifpoliveira@fei.edu.br, unifptakahashi@fei.edu.br, unievvianna@fei.edu.br}\par\smallskip
\orientador{Orientador: Charles Henrique Porto Ferreira}\par
\filiacao{Departamento de Ciência da Computação, Centro Universitário FEI}\par
\email{cferreira@fei.edu.br}\par

\resumo{Este trabalho investiga a validação intermediária de tarefas na adaptação do \textit{Ralph Wiggum Loop} para a resolução de perguntas em linguagem natural. O \textit{harness} proposto decompõe o enunciado em tarefas e utiliza o próprio modelo para avaliar cada resultado segundo critérios de aceite. Em 743 questões de múltipla escolha, a inferência orquestrada obteve 637 acertos, contra 652 na inferência direta. Nos 711 pares sem falha, o fluxo recuperou 35 respostas incorretas e perdeu 26 acertos da inferência direta. A análise de uma questão do ITA mostra como um critério introduzido no planejamento pode levar à reprovação sucessiva da alternativa correta.}
\keyword{modelos de linguagem, orquestração de inferência, critérios de aceite, validação iterativa}
\end{@twocolumnfalse}
]

\section{Introdução}

Em uma questão de matemática do ITA de 2024, o fluxo adaptado identificou a alternativa correspondente ao gabarito durante a execução de uma tarefa. A resposta final, entretanto, não chegou a ser produzida. Para compreender esse resultado, é necessário considerar como os critérios de aceite foram definidos a partir do enunciado.

A questão 43 da primeira fase solicitava um intervalo ao qual o raio de um cilindro \emph{poderia} pertencer. O modelo deduziu que $r>2$ e selecionou a alternativa E, que corresponde a \mbox{$2<r<2{,}5$}. Na etapa de Planejamento, o modelo havia estabelecido que a alternativa deveria abranger todos os valores admissíveis de $r$. Sob esse critério, o verificador reprovou a seleção até esgotar as cinco tentativas permitidas. A execução encerrou-se após 68,68 segundos. Em uma chamada direta, o mesmo modelo selecionou E em 2,81 segundos.

No \textit{Ralph Wiggum Loop} original, os requisitos do projeto e os testes de \textit{software} orientam a verificação das tarefas~\cite{huntley}. Na adaptação para perguntas em linguagem natural, o modelo formula os critérios que serão utilizados para verificar sua própria resolução. O objetivo experimental é avaliar como esse mecanismo afeta a acurácia e a capacidade do sistema de produzir uma resposta final.

\newpage
\section{Metodologia}

O \textit{harness}, camada responsável pela orquestração das chamadas ao modelo, organiza a inferência em três fases. No Planejamento, o modelo identifica os critérios de aceite gerais e decompõe a pergunta em tarefas com critérios individuais. Uma chamada subsequente verifica a cobertura do plano. Após sua aprovação, o sistema registra dados do enunciado, como valores, unidades e alternativas.

Na Execução em Loop, o modelo resolve cada tarefa e, em uma nova janela de contexto, verifica a validade técnica da solução e o atendimento aos critérios individuais. Uma reprovação devolve à tarefa o parecer do verificador. Uma aprovação permite armazenar o resultado e incorporar os aprendizados pontuais à memória das etapas seguintes. Na Criação da Resposta Final, a verificação passa a considerar o conjunto dos resultados. Caso sejam identificadas lacunas, o fluxo admite tarefas adicionais antes da síntese.

\section{Metodologia de avaliação}

Para cada questão, foram comparadas uma chamada direta ao DeepSeek-V4-Flash-0731 e uma execução mediada pelo \textit{harness}, com o mesmo enunciado e as mesmas alternativas. O gabarito foi utilizado apenas na avaliação da resposta final. Das 753 questões previstas, 743 apresentaram registros pareados, sendo 268 de provas brasileiras, 285 das 57 disciplinas do MMLU~\cite{mmlu} e 190 do GPQA Diamond~\cite{gpqa}.

A análise das falhas associou a alternativa registrada ao histórico das 106 questões em que o fluxo adaptado não obteve acerto. Para cada caso, foram examinados o critério aplicado, o parecer de verificação e a etapa de encerramento.

\clearpage
\section{Resultados da validação iterativa}

A acurácia nos 743 pares foi de 87,75\% na inferência direta e 85,73\% na inferência orquestrada. A Tabela~\ref{tab:pares} discrimina os acertos exclusivos de cada configuração e os acertos em comum. Considera-se válido o par em que ambas produziram uma alternativa passível de comparação com o gabarito. Na avaliação de ponta a ponta, uma falha operacional conta como erro para a configuração em que ocorreu.

\begin{table}[H]
\caption{Distribuição dos pares de questões segundo o acerto da inferência direta e da orquestrada.}\label{tab:pares}
\centering\small
\begin{tabular}{@{}cccc@{}}
\toprule
& \multicolumn{2}{c}{Pares} & \\
\cmidrule(lr){2-3}
Configuração que acertou & Sem falha & Com falha & Total \\
\midrule
Ambos & 601 & 0 & 601 \\
Somente orquestrada & 35 & 1 & 36 \\
Somente direta & 26 & 25 & 51 \\
Nenhum & 49 & 6 & 55 \\
\midrule
Total & 711 & 32 & 743 \\
\bottomrule
\end{tabular}
\end{table}

Nos 711 pares sem falha, o fluxo adaptado atingiu 89,45\% de acurácia, contra 88,19\% da inferência direta, com saldo de nove acertos. Nos 32 pares restantes, a inferência direta acertou 25 questões e a orquestrada, uma. A diferença de 24 acertos nesse segundo grupo explica a inversão do saldo na avaliação de ponta a ponta. O teste exato bilateral de McNemar, aplicado às questões em que apenas uma configuração acertou, não indicou diferença estatisticamente significativa a 5\% no conjunto completo ($p=0{,}1329$) ou no recorte sem falha ($p=0{,}3057$).

A Figura~\ref{fig:falhas} apresenta as causas de encerramento das 31 execuções em que o \textit{harness} falhou. A inferência direta registrou duas falhas, uma delas na mesma questão. O esgotamento das cinco tentativas de resolução de uma tarefa respondeu por 16 casos, dos quais 13 foram resolvidos corretamente pela inferência direta.

Em oito execuções, a saída intermediária não atendeu ao formato exigido. Em uma questão de biologia do ENEM de 2022, por exemplo, o verificador aprovou a alternativa correta, mas uma chave excedente no texto estruturado impediu o sistema de interpretar o parecer. Em outros dois casos, a execução chegou ao fim e a falha ocorreu na extração da alternativa a partir da resposta final.

\newpage

\begin{figure}[H]
\centering
\begin{tikzpicture}[x=0.20cm,y=0.53cm,font=\sffamily\footnotesize]
  \foreach \linha/\quantidade/\motivo in {
    5/16/{Tentativas de tarefa esgotadas},
    4/8/{Saída estruturada inválida},
    3/3/{Tempo de sessão excedido},
    2/2/{Resposta final inválida},
    1/1/{Tentativas de plano esgotadas},
    0/1/{Crédito do provedor esgotado}} {
    \node[anchor=east,inner sep=2pt] at (0,\linha) {\motivo};
    \fill[azulFEI!85!black] (0,\linha-.18) rectangle (\quantidade,\linha+.18);
    \node[anchor=west,inner sep=2pt] at (\quantidade,\linha) {\quantidade};
  }
\end{tikzpicture}
\caption{Causa terminal das 31 falhas operacionais registradas na inferência orquestrada.}\label{fig:falhas}
\end{figure}

\section{Discussão}

A aprovação das tarefas também ocorreu em trajetórias que produziram respostas incorretas. Das 26 questões sem falha em que apenas a inferência direta acertou, 20 foram concluídas sem qualquer reprovação de tarefa. Entre os 49 pares em que ambas as configurações produziram alternativas válidas e divergentes do gabarito, houve coincidência na alternativa incorreta em 41 casos. A reinicialização do contexto permite reexaminar cada resultado, mas a verificação continua a depender da interpretação do mesmo modelo.

Os históricos analisados em detalhe correspondem aos 106 casos de erro ou falha da inferência orquestrada. Essa seleção permite identificar mecanismos de falha, mas limita a estimativa de sua frequência entre as respostas corretas. A avaliação utilizou um único modelo e questões com gabarito. Nos pares sem falha, o consumo total de \textit{tokens} do \textit{harness} correspondeu a 98,36 vezes o da inferência direta.

\section{Conclusão}

A análise do caso do ITA evidencia uma limitação da revisão local. O verificador aplicou o critério definido no Planejamento, embora esse critério alterasse a condição solicitada no enunciado. Como continuidade, propõe-se avaliar o retorno ao Planejamento quando uma mesma objeção persiste entre tentativas. A comparação deverá considerar tanto as questões recuperadas pela revisão do critério quanto aquelas em que a alteração compromete um acerto anterior.

\section{Referências}
\printbibliography[heading=none]
\end{document}
